\documentclass[11pt]{beamer}


% ===============================
% Chinese + fonts (XeLaTeX/LuaLaTeX)
% ===============================
\usepackage[UTF8,scheme=plain,fontset=none]{ctex} % it handles fontspec loading for Chinese
\setmainfont{TeX Gyre Termes} % like Times New Roman \textrm{...}
\setsansfont{TeX Gyre Heros}  % like Arial \textsf{...}
\setmonofont{TeX Gyre Cursor} % like Consolas \texttt{...}

% CJK fonts (Noto fonts from Google)
\setCJKmainfont{Noto Serif CJK SC}  % Song Ti
\setCJKsansfont{Noto Sans CJK SC}   % Hei Ti \textsf{...}
\setCJKmonofont{Noto Sans Mono CJK SC} % Mono Font \texttt{...}


% ===============================
% page layout
% ===============================
\setbeamertemplate{navigation symbols}{}


% ===============================
% math packages
% ===============================
\usepackage{amsmath,bm}


% ===============================
% graphics
% ===============================
\usepackage{graphicx}
\usepackage{subcaption}
\graphicspath{{Images/}{./}}


% ===============================
% TikZ (for diagrams)
% ===============================
\usepackage{tikz}
\usetikzlibrary{
    arrows.meta,
    calc,
    positioning,
    decorations.pathmorphing,
    patterns,
    shapes.geometric,
    fit
}


% ===============================
% tables
% ===============================
\usepackage{booktabs}
\usepackage{array}
\usepackage{multirow}


% ===============================
% bibliography (biblatex)
% ===============================
\usepackage[backend=biber,style=authoryear,sorting=nyt]{biblatex}
\addbibresource{references.bib}


% ===============================
% hyperlinks
% (beamer loads hyperref automatically)
% ===============================
% \hypersetup{
%     colorlinks=true,
%     linkcolor=blue,
%     citecolor=blue,
%     urlcolor=blue
% }


% ===============================
% settings from latextemplates.com
% ===============================
\usetheme{Madrid} % layout theme
\usefonttheme{default}
\useinnertheme{circles}


% ===============================
% custom math commands (optional)
% ===============================
% \newcommand ---- define new command
% ===============================
% Basic sets of numbers
% ===============================
\newcommand{\R}{\mathbb{R}}
\newcommand{\N}{\mathbb{N}}
\newcommand{\Z}{\mathbb{Z}}
\newcommand{\Q}{\mathbb{Q}}

% ===============================
% Differential operators and constants
% ===============================
\newcommand{\dd}{\,\mathrm{d}}
\newcommand{\ee}{\mathrm{e}}
\newcommand{\ii}{\mathrm{i}}

% ===============================
% Vector calculus operators
% ===============================
\newcommand{\grad}{\nabla}
\newcommand{\diver}{\nabla\cdot}
\newcommand{\curl}{\nabla\times}


% ===============================
% Function spaces
% ===============================
\newcommand{\Hone}{H^1(\Omega)}
\newcommand{\Ltwo}{L^2(\Omega)}
\newcommand{\Hdiv}{H(\text{div},\Omega)}
\newcommand{\Hcurl}{H(\text{curl},\Omega)}


% ===============================
% Common operators (amsmath required)
% Put this in mathcmd.tex, and ensure amsmath is loaded before % \newcommand ---- define new command
% ===============================
% Basic sets of numbers
% ===============================
\newcommand{\R}{\mathbb{R}}
\newcommand{\N}{\mathbb{N}}
\newcommand{\Z}{\mathbb{Z}}
\newcommand{\Q}{\mathbb{Q}}

% ===============================
% Differential operators and constants
% ===============================
\newcommand{\dd}{\,\mathrm{d}}
\newcommand{\ee}{\mathrm{e}}
\newcommand{\ii}{\mathrm{i}}

% ===============================
% Vector calculus operators
% ===============================
\newcommand{\grad}{\nabla}
\newcommand{\diver}{\nabla\cdot}
\newcommand{\curl}{\nabla\times}


% ===============================
% Function spaces
% ===============================
\newcommand{\Hone}{H^1(\Omega)}
\newcommand{\Ltwo}{L^2(\Omega)}
\newcommand{\Hdiv}{H(\text{div},\Omega)}
\newcommand{\Hcurl}{H(\text{curl},\Omega)}


% ===============================
% Common operators (amsmath required)
% Put this in mathcmd.tex, and ensure amsmath is loaded before \input{mathcmd}
% ===============================
% --- Linear algebra / matrix analysis ---
\DeclareMathOperator{\rank}{rank}
\DeclareMathOperator{\tr}{tr}
\DeclareMathOperator{\diag}{diag}
\DeclareMathOperator{\trace}{trace}


% ===============================
% --- Linear algebra / matrix analysis ---
\DeclareMathOperator{\rank}{rank}
\DeclareMathOperator{\tr}{tr}
\DeclareMathOperator{\diag}{diag}
\DeclareMathOperator{\trace}{trace}




\title[Short Title]{Full Presentation Title}
\subtitle{Optional Subtitle}
\author[James Cook \and Roald Amundsen]{Capt. James Cook \and Roald Amundsen}
\institute[UC]{University of Cambridge \\ \smallskip \textit{james@LaTeXTemplates.com}}
\date[\today]{International Symposium of Explorers \\ \today}

% ============================
\begin{document}
% ============================



%------------------------------------------------
\begin{frame}
	\titlepage
\end{frame}
%------------------------------------------------



%------------------------------------------------
\begin{frame}
	\frametitle{Presentation Overview}
	\tableofcontents
\end{frame}
%------------------------------------------------



\section{Text Examples}
\subsection{Paragraphs and Lists}
%------------------------------------------------
\begin{frame}{中文测试}
	\begin{quote}
		Sed diam enim, sagittis nec \alert{condimentum} sit amet, ullamcorper sit amet libero. Aliquam vel dui orci, a porta odio.\\
		--- Someone, somewhere\ldots
	\end{quote}
	\bigskip % Vertical whitespace
	Nullam id suscipit ipsum. Aenean lobortis commodo sem, ut commodo leo gravida vitae. Pellentesque vehicula ante iaculis arcu pretium rutrum eget sit amet purus. Integer ornare nulla quis neque ultrices lobortis.
\end{frame}
%------------------------------------------------



\subsection{Columns}
%------------------------------------------------
\begin{frame}{Multiple Columns}
	\begin{columns}[c] % The "c" option specifies centered vertical alignment while the "t" option is used for top vertical alignment
		\begin{column}{0.45\textwidth} % Left column width
			\textbf{Heading}
			\begin{enumerate}
				\item Statement
				\item Explanation
				\item Example
			\end{enumerate}
		\end{column}
		\begin{column}{0.5\textwidth} % Right column width
			Lorem ipsum dolor sit amet, consectetur adipiscing elit. Integer lectus nisl, ultricies in feugiat rutrum, porttitor sit amet augue. Aliquam ut tortor mauris. Sed volutpat ante purus, quis accumsan dolor.
		\end{column}
	\end{columns}
\end{frame}
%------------------------------------------------



\section{Table and Figure Examples}
%------------------------------------------------
\begin{frame}
    \frametitle{Table \& Citation Example}
	\framesubtitle{Subtitle} % Optional subtitle
	\begin{table}
		\begin{tabular}{l l l}
			\toprule
			\textbf{Treatments} & \textbf{Response 1} & \textbf{Response 2}\\
			\midrule
			Treatment 1 & 0.0003262 & 0.562 \\
			Treatment 2 & 0.0015681 & 0.910 \\
			Treatment 3 & 0.0009271 & 0.296 \\
			\bottomrule
		\end{tabular}
		\caption{Table caption}
	\end{table}
    An example of the \texttt{\textbackslash textcite} or \texttt{\textbackslash parencite} command to cite within the presentation and \texttt{\textbackslash footfullcite} to show the bibliography on the bottom: This statement requires citation \parencite{p1,p2}\footfullcite{p1}\footfullcite{p2}
\end{frame}
%------------------------------------------------



\section{Mathematics}
\subsection{Blocks}
%------------------------------------------------
\begin{frame}{Blocks of Highlighted Text \& Definitions}
	\begin{block}{Block Title} % first parameter: block, exampleblock, alertblock; 
        % second parameter: title(can be empty)
		Lorem ipsum dolor sit amet, consectetur adipiscing elit. Integer lectus nisl, ultricies in feugiat rutrum, porttitor sit amet augue.
	\end{block}
	\begin{block}{} % Block without title
		Suspendisse tincidunt sagittis gravida. Curabitur condimentum, enim sed venenatis rutrum, ipsum neque consectetur orci.
	\end{block}
	\begin{definition}
		A \alert{prime number} is a number that has exactly two divisors.
	\end{definition}
	% \smallskip % Vertical whitespace
	You can also use the \texttt{theorem}, \texttt{lemma}, \texttt{proof} and \texttt{corollary} environments.
\end{frame}
%------------------------------------------------



%------------------------------------------------
\begin{frame}[allowframebreaks]{References}
    % Use [allowframebreaks] to allow automatic splitting across slides if the content is too long
    \footnotesize
    \printbibliography
\end{frame}
%------------------------------------------------



%------------------------------------------------
\begin{frame}{test url}
    \begin{itemize}
        \item \url{https://www.latextemplates.com}
        \item \href{https://www.latextemplates.com}{LaTeX Templates}
    \end{itemize}
\end{frame}
%------------------------------------------------



%------------------------------------------------
\begin{frame}[plain] % The optional argument 'plain' hides the headline and footline
	\begin{center}
        {\Huge Thank You!}

		\bigskip\bigskip % Vertical whitespace

		{\LARGE Questions? Comments?}
	\end{center}
\end{frame}
%------------------------------------------------



% ============================
\end{document} 
% ============================
